\section{Open Questions: cuándo recuperar, cómo escalar y si la evidencia se usa}

Después de Atlas, la clase pasa a las preguntas abiertas, intercaladas con preguntas y respuestas de los estudiantes. Aquí el foco deja de ser «qué arquitectura es más fuerte» y pasa a la policy dinámica, el scaling, la legal-risk motivation, la posición en el context, el tool use y el instruction tuning. Muchas de estas preguntas todavía no pueden responderse con un solo benchmark.

\subsection{RAG y Long Context no son del todo opuestos}

Un estudiante pregunta si un context muy largo reemplazará a RAG. El ponente responde que meter todo Harry Potter en el context solo para encontrar el nombre de una lechuza es muy ineficiente; además, para escalar la attention, los métodos de contexto largo suelen introducir sparsity, lo que en esencia también equivale a seleccionar unas pocas conexiones relevantes.

\begin{knowledgebox}{El problema común de Long context y retrieval}
Ambos tienen que decidir «qué tokens vale la pena calcular». Long context coloca esa selección en el attention pattern o en la gestión de la cache; RAG la coloca en un external index. Al final, los dos enfrentan relevance, latency, position bias y evidence verification.
\end{knowledgebox}

\teachervoice{La conclusión de la clase es que, si el contexto largo llega a escalar a tamaños enormes, se parecerá cada vez más a la sparse/non-parametric retrieval, en lugar de aplicar full attention a todos los tokens sin costo.}

\subsection{When to Retrieve: convertir la recuperación en una estrategia}

Recuperar en cada token desperdicia compute, mientras que recuperar una sola vez al principio puede pasar por alto necesidades de conocimiento posteriores. FLARE permite que el modelo decida activamente cuándo buscar y qué buscar, según la confianza de la generación o el contenido que viene, de modo que la retrieval deja de ser un pipeline fijo y se vuelve una sequential policy.

\lecturefigure{slide-39-when-to-retrieve.jpg}{FLARE y active retrieval: decidir cuándo y por qué recuperar}{Video de la clase de Stanford 00:56:36}

\begin{importantbox}{El objetivo de costo de la retrieval policy}
Sea la acción $a_t\in\{\text{retrieve},\text{continue}\}$, el costo de recuperación $c_R$ y el costo de error $c_E$; puede optimizarse
\[
\mathbb{E}\!\left[\sum_t c_R\mathbf{1}(a_t=\text{retrieve})
+c_E\mathbf{1}(\hat y\ne y)\right].
\]
La policy debe aprender a invocar la herramienta cuando el valor de la información supera el costo de recuperarla, en lugar de suponer que más retrieval calls son mejores por naturaleza.
\end{importantbox}

\subsection{How to Train at Scale}

La Scale slide analiza la full async update, la query-only update, TRIME y métodos aproximados, así como el entrenamiento in-batch tras construir lotes similares con BM25. Cuanto más grande es el index, más cara resulta la full refresh del lado de los documentos; por eso el hard-negative sampling, el memory bank y el reranker se vuelven aproximaciones importantes.

\lecturefigure{slide-40-training-at-scale.jpg}{Métodos in-batch y de memoria para retriever training a gran escala}{Video de la clase de Stanford 00:57:39}

\begin{knowledgebox}{Lectura de la figura: qué ahorran los in-batch negatives}
Los demás documents de un lote pueden servir como negatives, lo que evita buscar por separado una gran cantidad de ejemplos negativos para cada query; el memory bank amplía el conjunto de negativos. El riesgo son los false negatives: documentos semánticamente relevantes tratados como negativos, lo cual distorsiona la representation.
\end{knowledgebox}

\subsection{SILO: llevar la legal-risk motivation a la arquitectura}

SILO propone entrenar el parametric model con datos más seguros y de licencia más clara, y luego, en tiempo de prueba, acceder a contenido más amplio desde un non-parametric datastore. La clase lo considera una dirección elegante para atender la copyright provenance y el legal risk, pero no afirma que el contenido del index quede por ello libre de responsabilidad legal.

\lecturefigure{slide-41-silo-legal-risk.jpg}{SILO: aislar parte del riesgo de los datos de entrenamiento con un non-parametric datastore}{Video de la clase de Stanford 00:58:45}

\begin{warningbox}{La Architecture no resuelve automáticamente los problemas legales}
El datastore todavía debe considerar autorización, acceso, registros, eliminación, citas en la salida y datos de usuarios; además, el modelo puede reescribir el retrieved content. El valor de investigación de SILO está en separar los datos de entrenamiento del test-time knowledge, lo que facilita reemplazarlos y auditarlos; no es una conclusión legal.
\end{warningbox}

\teachervoice{El ponente relaciona SILO con las demandas de aquel momento sobre el origen de los datos de entrenamiento de los modelos, y enfatiza que el parametric model puede entrenarse con datos public-domain y luego recuperar de un index de mayor riesgo. Estas notas conservan esa motivation, pero dejan claro que no la elevan a garantía de cumplimiento normativo.}

\subsection{Lost in the Middle: aunque la recuperación sea correcta, el uso puede fallar}

La curva de la Lost-in-the-Middle slide muestra que el relevant document se aprovecha mejor cuando está al principio o al final del context, y que el desempeño baja cuando está en medio. Si el generator respetara por completo el ranking del retriever, el cambio de posición no tendría un efecto tan marcado; esto demuestra que la evidence placement y el attention behavior son un failure mode independiente.

\lecturefigure{slide-42-lost-in-middle.jpg}{Un relevant passage en medio de un context largo tiende a ignorarse}{Video de la clase de Stanford 01:00:45}

\begin{knowledgebox}{Lectura de la figura: qué indica la U-shaped curve}
El eje horizontal es la posición del relevant document entre 20 retrieved documents; el eje vertical, el desempeño en la tarea. Valores altos al principio y al final y bajos en medio indican que el modelo tiene position bias. Esto no demuestra que los documentos intermedios nunca sirvan, sino que el reranking, el context packing y la evaluation deben considerar la posición al mismo tiempo.
\end{knowledgebox}

\begin{warningbox}{Retrieval recall y answer grounding deben medirse por capas}
Puede registrarse si la evidence clave entra en el top-$k$, si llega al prompt final, en qué posición queda, si la respuesta la cita y si la salida es coherente con la evidence. Mirar solo el final exact match no permite ubicar si la falla está en search, en packing o en generation.
\end{warningbox}

\subsection{Toolformer: la Retrieval es un Tool Use más general}

La Toolformer slide coloca la search API dentro de un conjunto de herramientas como calculator, calendar y translation. El Language model puede aprender cuándo insertar una API call, leer el valor devuelto y continuar generando; así, la retrieval es una instancia de tool augmentation.

\lecturefigure{slide-43-toolformer.jpg}{Toolformer: un modelo de lenguaje aprende de forma autosupervisada a invocar herramientas externas}{Video de la clase de Stanford 01:01:57}

\begin{knowledgebox}{Lectura de la figura: el cambio de interfaz de RAG a Agent}
El RAG fijo invoca la recuperación en cada solicitud; un tool-using LM puede elegir search, calculator u otra API. El sistema incorpora action selection, argument generation, tool error handling y permission boundary, y el problema pasa de una architecture estática a la sequential decision making.
\end{knowledgebox}

\teachervoice{El ponente señala que la active web search no necesariamente consulta un index propio; de manera más general, el LM es un tool user, la retrieval es solo una de muchas herramientas y aprender de verdad la selección de acciones podría requerir RL.}

\subsection{Self-RAG: Retrieve, Generate, Critique}

Self-RAG hace que un único LM genere reflection tokens para decidir si recupera, evaluar la evidence relevance, generar la respuesta y hacer la critique de support/utility. Busca integrar la active retrieval y la autoevaluación en un mismo marco de entrenamiento. Esta sección retoma la selección de acciones de Toolformer y se pregunta además cómo un mismo modelo juzga la calidad de la evidencia y el grado de sustento de la respuesta.

\lecturefigure{slide-44-self-rag.jpg}{El flujo de recuperación activa y reflection de Self-RAG}{Video de la clase de Stanford 01:02:15}

\begin{importantbox}{La Self-reflection sigue necesitando evaluación externa}
El Critique token es una predicción del modelo, no un oráculo de hechos. El sistema debe verificar por separado la retrieval decision, la source relevance, el claim support y la final utility; la autoevaluación del modelo puede usarse como feature, pero no sustituye a las held-out labels, la human review ni un verifier ejecutable.
\end{importantbox}

\begin{lstlisting}[caption={Ciclo de active retrieval y evidence verification}]
state = initialize(question)
while not state.finished and state.steps < max_steps:
    if policy.needs_evidence(state):
        docs = search(state.query)
        docs = rerank_and_filter(docs, source_policy)
        state.add_evidence(docs)
    draft = generator.continue_answer(state)
    verdict = verify_claims(draft, state.evidence)
    state = revise_or_finish(state, draft, verdict)
return state.answer
\end{lstlisting}

\subsection{Instruction Tuning the Entire RAG System}

La Instruction-tuning slide compara InstructRetro con RA-DIT: antes, los datos de instrucciones entrenaban sobre todo al generator, y el retriever no necesariamente entendía la intención de la tarea, como «comparar, citar, usar solo las fuentes indicadas». El Dual instruction tuning adapta a la vez la retrieval y la generation, para que todo el sistema siga la instruction.

\lecturefigure{slide-45-instruction-tuning.jpg}{InstructRetro y RA-DIT: instruction tuning del RAG system}{Video de la clase de Stanford 01:03:03}

\begin{knowledgebox}{La Instruction también tiene semántica para el Retriever}
``Give a counterargument'', ``use peer-reviewed sources'' y ``summarize this document'' requieren evidence distinta. Si el retriever solo atiende a la topic similarity, puede devolver documentos relacionados con el tema que no cumplen el task operator; una instruction-aware query representation puede reducir ese desajuste.
\end{knowledgebox}

\subsection{Advanced Frozen RAG: innovación práctica de la Developer Community}

Aun sin joint training, los desarrolladores han construido el child-parent recursive retriever, la hybrid search/RRF, el zero-shot LLM reranker y HyDE. HyDE primero hace que el modelo genere un hypothetical document y luego usa su embedding para buscar documentos reales; la clase lo llama una idea interesante de «usar la hallucination para corregir la hallucination».

\lecturefigure{slide-46-advanced-frozen-rag.jpg}{Recuperación jerárquica, fusión, reranking y HyDE en Advanced Frozen RAG}{Video de la clase de Stanford 01:04:30}

\begin{knowledgebox}{Términos clave: cuatro advanced frozen patterns}
\begin{tabularx}{\textwidth}{L{0.25\textwidth}>{\raggedright\arraybackslash}X}
\toprule
Patrón & Mecanismo \\
\midrule
Child-parent retrieval & Recupera con precisión mediante chunks pequeños y luego se amplía al párrafo padre para aportar el context completo \\
Hybrid/RRF & Combina el ranking de BM25 y el dense ranking, para cubrir tanto el exact match como el semantic match \\
LLM reranker & Hace que un LM potente compare la relevance dentro de un conjunto pequeño de candidatos; su costo es alto \\
HyDE & Primero genera una hypothetical answer/document y luego recupera evidencia real con su embedding \\
\bottomrule
\end{tabularx}
\end{knowledgebox}

\begin{warningbox}{El hypothetical text de HyDE no es evidencia}
Solo sirve para construir una query representation más rica; la respuesta final debe seguir citando retrieved documents reales. Si el hypothetical document se toma directamente como un hecho, la retrieval query expansion se convierte en un ciclo autoconfirmatorio.
\end{warningbox}

\teachervoice{Por un lado, el ponente critica el carácter de Frankenstein del frozen RAG; por otro, reconoce la innovación práctica del ecosistema de desarrolladores, como LlamaIndex, LangChain, Chroma y Weaviate. La dirección de investigación y el baseline de ingeniería pueden ser válidos al mismo tiempo.}

\subsection{Resumen de la sección}

Las Open questions se concentran en la policy y en los límites del sistema: cuándo recuperar, cómo escalar el retriever, cómo manejar la legal-risk motivation, cómo asegurar que el generator use de verdad la evidencia y cómo la retrieval se convierte en tool use. El Frozen RAG todavía cuenta con técnicas de ingeniería sólidas, pero el joint learning y la evaluation por capas determinan su techo.
